% AUTO-GENERATED. Edit the Markdown sources or notes.config.json.
% Sources: 01-Mathematics/2026-Probability-Theory/2026-09-23-Note.md

\section{随机变量的变换与期望}

\subsection{Distributions of Functions of a Random Variable}

设随机变量 \(X\) 的分布已知，并定义

\[
Y=g(X).
\]

目标是由 \(X\) 的分布求出 \(Y\) 的分布。首先要确定 \(g\) 在 \(X\) 的 support 上是否 one-to-one，再选择离散求和、CDF method 或 change of variables。

\subsubsection{Discrete Transformations}

\textbf{Theorem}

\begin{quote}
If \(X\) is discrete and \(Y=g(X)\), then
\[
p_Y(y)=\sum_{x:g(x)=y}p_X(x).
\]
\end{quote}

同一个 \(y\) 可能对应多个 \(x\)。因此，\(Y=y\) 的概率等于所有满足 \(g(x)=y\) 的原像点概率之和。

\textbf{Example}

设 \(X\) 在集合 \(\{-2,-1,0,1,2\}\) 上均匀分布，令 \(Y=X^2\)。则

\[
\begin{aligned}
P(Y=0)&=P(X=0)=\frac15,\\
P(Y=1)&=P(X=-1)+P(X=1)=\frac25,\\
P(Y=4)&=P(X=-2)+P(X=2)=\frac25.
\end{aligned}
\]

\(g(x)=x^2\) 不是 one-to-one，所以除 \(0\) 外，每个 \(Y\) 的可能取值都要合并两个原像点的概率。

\subsubsection{CDF Method}

\textbf{Definition}

\begin{quote}
For \(Y=g(X)\), begin with
\[
F_Y(y)=P(g(X)\le y).
\]
Rewrite the event in terms of \(X\), then differentiate when a density exists.
\end{quote}

CDF method 直接处理事件 \(\{g(X)\le y\}\)，不要求 \(g\) 单调。它特别适用于非单调变换、最大值或最小值，以及变换后 support 随 \(y\) 改变的情形。

求解时需要先把不等式改写为关于 \(X\) 的事件：

\[
\begin{aligned}
F_Y(y)
&=P(Y\le y)\\
&=P(g(X)\le y).
\end{aligned}
\]

若最后得到的 \(F_Y\) 可微，则

\[
f_Y(y)=F_Y'(y).
\]

\subsubsection{One-to-One Change of Variables}

\textbf{Theorem}

\begin{quote}
Suppose \(Y=g(X)\) and \(g\) is strictly monotone and differentiable. If \(x=g^{-1}(y)\), then
\[
f_Y(y)=f_X\bigl(g^{-1}(y)\bigr)
\left|\frac{d}{dy}g^{-1}(y)\right|
\]
for \(y\) in the transformed support.
\end{quote}

绝对导数

\[
\left|\frac{d}{dy}g^{-1}(y)\right|
\]

是 one-dimensional Jacobian。它修正了变量变换对区间长度的伸缩，使变换前后的概率保持一致。

\textbf{Example}

设 \(U\sim\operatorname{Unif}(0,1)\)，并令

\[
Y=-\log U.
\]

由 \(u=e^{-y}\) 可知 \(Y\) 的 support 为 \(y>0\)，且

\[
\left|\frac{du}{dy}\right|=e^{-y}.
\]

因此

\[
\begin{aligned}
f_Y(y)
&=f_U(e^{-y})\left|\frac{du}{dy}\right|\\
&=e^{-y},\qquad y>0.
\end{aligned}
\]

故

\[
Y\sim\operatorname{Exp}(1).
\]

也可用 CDF method 验证：对 \(y>0\)，

\[
\begin{aligned}
F_Y(y)
&=P(-\log U\le y)\\
&=P(U\ge e^{-y})\\
&=1-e^{-y}.
\end{aligned}
\]

\subsubsection{Non-One-to-One Transformations}

\textbf{Theorem}

\begin{quote}
If \(g\) has inverse branches \(x_j(y)\), then
\[
f_Y(y)=\sum_j f_X\bigl(x_j(y)\bigr)
\left|\frac{dx_j}{dy}\right|.
\]
\end{quote}

当一个 \(y\) 对应多个原像点时，每个 inverse branch 都贡献一项，最后把这些密度相加。

\textbf{Example}

令 \(Y=X^2\)。对 \(y>0\)，两个 inverse branches 为

\[
x_1(y)=\sqrt y,
\qquad
x_2(y)=-\sqrt y.
\]

相应的 Jacobian 绝对值均为 \(1/(2\sqrt y)\)，所以

\[
f_Y(y)
=\frac{f_X(\sqrt y)+f_X(-\sqrt y)}{2\sqrt y},
\qquad y>0.
\]

\subsubsection{Normal and Chi-Square}

\textbf{Theorem}

\begin{quote}
If \(Z\sim N(0,1)\), then
\[
Z^2\sim\chi_1^2.
\]
\end{quote}

令 \(Y=Z^2\)。标准正态密度关于原点对称，因此对 \(y>0\)，

\[
\begin{aligned}
f_Y(y)
&=\frac{\phi(\sqrt y)+\phi(-\sqrt y)}{2\sqrt y}\\
&=\frac{1}{\sqrt{2\pi y}}e^{-y/2}.
\end{aligned}
\]

这正是自由度为 \(1\) 的 chi-square density。

进行变量变换时，应先写出 \(X\) 的 support，再求 \(Y=g(X)\) 的 support，并列出全部 inverse branches。所得 PMF 或 PDF 还应满足非负性和总质量为 \(1\)。

\subsection{Expected Values}

\subsubsection{LOTUS}

\textbf{Theorem}

\begin{quote}
For a function \(g\) for which the expectation exists,
\[
E[g(X)]=\sum_x g(x)p_X(x)
\]
for discrete \(X\), and
\[
E[g(X)]=\int_{-\infty}^{\infty}g(x)f_X(x)\,dx
\]
for continuous \(X\).
\end{quote}

该结论称为 \textbf{Law of the Unconscious Statistician (LOTUS)}。计算 \(E[g(X)]\) 时，可以直接对 \(X\) 的分布求和或积分，不必先推导 \(Y=g(X)\) 的分布。

例如，若 \(X\) 连续且二阶矩存在，则

\[
E[X^2]=\int_{-\infty}^{\infty}x^2f_X(x)\,dx.
\]

\subsubsection{Linearity of Expectation}

\textbf{Theorem}

\begin{quote}
Whenever the expectations exist,
\[
E[aX+bY+c]=aE[X]+bE[Y]+c.
\]
Independence is not required.
\end{quote}

更一般地，

\[
E\left[\sum_{i=1}^{n}a_iX_i\right]
=\sum_{i=1}^{n}a_iE[X_i].
\]

期望的线性性只要求相关期望存在，不要求随机变量之间独立。独立性会在乘积期望或方差加法中发挥作用，但不是线性性的前提。

\subsubsection{Indicator Variables}

\textbf{Definition}

\begin{quote}
For an event \(A\), let \(\mathbf 1_A\) equal \(1\) when \(A\) occurs and \(0\) otherwise. Then
\[
E[\mathbf 1_A]=P(A).
\]
\end{quote}

示性变量把事件是否发生转化为取值为 \(0\) 或 \(1\) 的随机变量。若

\[
N=\sum_{i=1}^{n}\mathbf 1_{\{X_i\in A\}}
\]

表示落入集合 \(A\) 的观测个数，则

\[
\begin{aligned}
E[N]
&=\sum_{i=1}^{n}E\left[\mathbf 1_{\{X_i\in A\}}\right]\\
&=\sum_{i=1}^{n}P(X_i\in A).
\end{aligned}
\]

若 \(X_1,\ldots,X_n\) identically distributed，且 \(P(X_i\in A)=p\)，则

\[
E[N]=np.
\]

这里仍不需要独立性；只要每一项的边缘概率为 \(p\) 即可。

\subsubsection{Variance}

\textbf{Definition}

\begin{quote}
The variance of \(X\) is
\[
\operatorname{Var}(X)
=E\left[(X-E[X])^2\right]
=E[X^2]-E[X]^2.
\]
\end{quote}

方差度量随机变量围绕均值的平方离散程度。由定义可得

\[
\operatorname{Var}(aX+b)=a^2\operatorname{Var}(X).
\]

平移常数 \(b\) 不改变离散程度，而缩放 \(a\) 会使方差乘以 \(a^2\)。

\subsubsection{Mean Minimizes Squared Distance}

\textbf{Theorem}

\begin{quote}
If \(E[X^2]<\infty\), then the constant \(a\) minimizing
\[
E[(X-a)^2]
\]
is \(a=E[X]\).
\end{quote}

\textbf{Proof}

令 \(\mu=E[X]\)，则

\[
\begin{aligned}
E[(X-a)^2]
&=E[(X-\mu+\mu-a)^2]\\
&=E[(X-\mu)^2]+2(\mu-a)E[X-\mu]+(\mu-a)^2\\
&=\operatorname{Var}(X)+(\mu-a)^2.
\end{aligned}
\]

第一项与 \(a\) 无关，第二项在且仅在 \(a=\mu\) 时为 \(0\)，所以均值是平方损失下的最优常数预测。
